\documentclass[aspectratio=169]{beamer}
\usepackage{tikz}
\usepackage{graphicx}
\usepackage{listings}
\usepackage{xcolor}
\usepackage{hyperref}
\usepackage[spanish,es-noquoting]{babel}
\usepackage[round]{natbib}

\lstset{
  basicstyle=\ttfamily\small,
  keywordstyle=\color{blue},
  stringstyle=\color{red},
  commentstyle=\color{gray},
  showstringspaces=false
}
% Minimal look
\usetheme{Warsaw}
%\setbeamertemplate{headline}
\setbeamertemplate{footline}{
  \leavevmode%
  \hbox{%
    % Left: Section name
    \begin{beamercolorbox}[wd=.33\paperwidth,ht=2.25ex,dp=1ex,left]{author in head/foot}%
      \hspace*{2ex}\insertsection
    \end{beamercolorbox}%
    % Center: Presentation title
    \begin{beamercolorbox}[wd=.34\paperwidth,ht=2.25ex,dp=1ex,center]{title in head/foot}%
      \insertshorttitle
    \end{beamercolorbox}%
    % Right: Slide numbers
    \begin{beamercolorbox}[wd=.33\paperwidth,ht=2.25ex,dp=1ex,right]{date in head/foot}%
      \insertframenumber{} / \inserttotalframenumber\hspace*{2ex}
    \end{beamercolorbox}%
  }%
  \vskip0pt%
}

\usecolortheme{default}
\usefonttheme{serif}
\graphicspath{{../}{./}}
\setbeamertemplate{caption}[numbered]

% Title info
\title[LegalAI Presentation]{LegalAI Presentation}
\subtitle{}
\author{Ricardo Reyes-Jimenez}
\institute{Universidad de La Sabana}
\date{\today}



\begin{document}

% Title slide
\begin{frame}
  \titlepage
\end{frame}

% Reset background for normal slides
\usebackgroundtemplate{}



\begin{frame}[fragile]{Research Brief: Análisis de datos y toma de decisiones usando IA para público en general}

This research brief compiles recent academic insights regarding the use of Artificial Intelligence (AI) for data analysis and decision-making, particularly concerning the general public. While direct local PDF context was not identified in the available library, a selection of high-quality academic papers provides a foundational understanding of current trends and challenges.

### Academic Citations

1.  \textbf{Aparicio Araque, Blanca. ``EL IMPACTO DE LA INTELIGENCIA ARTIFICIAL Y LA AUTOMATIZACIÓN DE DATOS EN LA TOMA DE DECISIONES EN EL SECTOR PÚBLICO.'' \textit{Inteligencia y toma de decisiones: perspectivas actuales}, 2025.}
\begin{itemize}
  \begin{itemize}
    \item \textbf{DOI:} 10.69592/979-13-7011-369-8-cap-17
    \item \textbf{Relevance:} This paper directly addresses the impact of AI and data automation on decision-making within the public sector, a realm that profoundly affects the general public. Its forward-looking publication year (2025) suggests it captures contemporary and emerging perspectives on the topic.
  \end{itemize}
\end{itemize}
2.  \textbf{Méndez Saldaña, M. ``Análisis de datos para la toma de decisiones en entornos virtuales.'' \textit{Unknown Journal}, N/A. DOI: 10.7764/tesisuc/edu/74749.}
\begin{itemize}
  \begin{itemize}
    \item \textbf{Relevance:} This academic work, likely a thesis, explores data analysis for decision-making in virtual environments. Given the increasing digitalization of public interaction and services, understanding how data analysis informs decisions in these virtual spaces is crucial for the general public's experience.
  \end{itemize}
\end{itemize}
3.  \textbf{Zouak Lara, P. ``Sinergias entre el RGPD y la IA en el ámbito sanitario para la toma de decisiones automatizadas.'' \textit{Inteligencia artificial y protección de datos: desafíos en la era digital}, 2025.}
\begin{itemize}
  \begin{itemize}
    \item \textbf{DOI:} 10.69592/979-13-7011-111-3-cap-16
    \item \textbf{Relevance:} While specifically focused on the healthcare sector and GDPR, this paper delves into the critical intersection of AI, automated decision-making, and data protection. These aspects are paramount for the general public as AI becomes more integrated into sensitive areas of their lives, highlighting the need for robust regulatory frameworks and ethical considerations.
  \end{itemize}
\end{itemize}
### Interesting Facts for Presentation

1.  \textbf{AI's Silent Decision-Makers:} Did you know that AI algorithms are increasingly making decisions that impact the general public's daily lives, from optimizing traffic flow and allocating public resources to personalizing online experiences, often without individuals explicitly realizing the AI's influence? This ``invisible hand'' of AI highlights the growing need for public literacy and ethical oversight in AI development.
2.  \textbf{The Transparency Challenge:} One of the biggest challenges in AI-driven public decision-making is ensuring transparency. Imagine an AI system denying a public service application; explaining \textit{why} it made that decision can be incredibly complex due to the intricate nature of machine learning models. This 'black box' problem raises significant questions about accountability and fairness when AI is used for the general public's benefit.

----------

\end{frame}



\begin{frame}[fragile]{}
title: ``La Magia Detrás de las Decisions: IA y Datos para Todos''
subtitle: ``Desentrañando el Futuro de la Toma de Decisiones en la Era Digital''
author: ``Prof. [Tu Nombre], Cátedra de Derecho y Tecnología''
date: ``Universidad del Mañana | [Fecha Actual]''
revealjs: true
\end{frame}



\begin{frame}[fragile]{¡Bienvenidos, Futuros Maestros de la IA!}


Hola a todos, ¿listos para un viaje fascinante?

Soy su profesor, y hoy no hablaremos de leyes aburridas ni de códigos crípticos. ¡Hoy desentrañaremos uno de los misterios más emocionantes de nuestro tiempo!

Piensen por un momento... ¿alguna vez han sentido que una fuerza invisible guía sus decisiones diarias? ¿Quizás su playlist favorita, la ruta de tráfico que eligen o incluso la serie que les recomienda su plataforma?

\textbf{¡Atención!} No es brujería... ¡es \textbf{Inteligencia Artificial y Datos}! Y está redefiniendo cómo se toman las decisiones, desde las más triviales hasta las que impactan a toda una ciudad.

\textbf{IMG\_PH\_0}

\end{frame}



\begin{frame}[fragile]{La IA Invisible: ¿Quién Decide por Ti?}


Imagina que te levantas un día y el semáforo en tu ruta usual ha cambiado su patrón. O que el hospital de tu zona ha optimizado sus citas de una forma increíblemente eficiente. O incluso, que tu gobierno local ha asignado recursos públicos de manera más justa que nunca.

\textbf{¿Magia? ¡No! Es la IA.}

La Inteligencia Artificial, sin que nos demos cuenta, se ha convertido en una especie de `\texttt{mano invisible'' que optimiza muchísimos aspectos de nuestra vida. Desde predecir qué autobús necesitas hasta cómo se distribuyen ciertos servicios públicos }\cite{Aparicio Araque, 2025}`.

\textbf{Pregunta Interactiva:} ¿Quién de aquí ha sentido que una app ``adivinó'' lo que quería hacer o comprar? ¡Levanten la mano!

\textbf{IMG\_PH\_0}

\end{frame}



\begin{frame}[fragile]{¿Qué es esta ``IA'' y qué son los ``Datos''? ¡Desmitificando!}


¡No se asusten con los términos! Piénsenlo así:

\begin{itemize}
  \item \textbf{Inteligencia Artificial (IA): El Cerebro Digital}
  \begin{itemize}
    \item Es como un estudiante súper aplicado y rapidísimo que aprende de ejemplos (datos). No tiene sentimientos, pero es una máquina de patrones.
    \item Le das miles de fotos de gatos, y aprenderá a reconocer un gato. Le das miles de rutas de tráfico y sabrá cuál es la más rápida.
    \item Su objetivo: resolver problemas y tomar decisiones.
  \end{itemize}
  \item \textbf{Análisis de Datos: El Trabajo de Detective}
  \begin{itemize}
    \item Los datos son las ``pistas''. Cada compra, cada clic, cada interacción, cada sensor de tráfico... ¡es un dato!
    \item El análisis de datos es el arte de examinar esas pistas para encontrar patrones, tendencias y extraer información valiosa.
    \item Es lo que permite a la IA aprender y mejorar. Sin datos, la IA está ciega.
  \end{itemize}
\end{itemize}
\textbf{En resumen:} Los datos son el \textit{combustible}, y la IA es el \textit{motor} que usa ese combustible para \textit{tomar decisiones}.

\textbf{IMG\_PH\_0}

\end{frame}



\begin{frame}[fragile]{El Poder Oculto de los Datos: Una Pequeña Historia}


Había una vez una pequeña ciudad que sufría de constantes atascos de tráfico. Los alcaldes probaban de todo: más semáforos, menos semáforos, policía... ¡nada funcionaba!

Hasta que un día, un joven ingeniero propuso algo radical: ``Recopilemos \textit{todos} los datos de tráfico. Cada coche, cada hora, cada ruta.''
Al principio, la gente se rió. ``¿Más datos? ¿Para qué?''

Pero el ingeniero usó un nuevo sistema (¡una forma temprana de IA!) para analizar billones de puntos de datos. Descubrió que el problema no eran los semáforos, sino un pequeño desvío que la gente tomaba para evitar una escuela, pero que congestionaba una calle principal en horas pico. Con un cambio mínimo en la señalización y un semáforo inteligente, ¡el tráfico fluyó!

\textbf{Moraleja:} Los datos tienen el poder de revelar verdades ocultas y ofrecer soluciones sorprendentes. ¡Son la clave para decisiones inteligentes!

\textbf{IMG\_PH\_0}

\end{frame}



\begin{frame}[fragile]{IA en Acción: Decisiones que nos Toca Vivir}


La IA ya está tomando decisiones que nos afectan directamente, a veces para bien, a veces con desafíos.

\begin{itemize}
  \item \textbf{Optimizando el Tráfico Urbano:} Analiza millones de datos en tiempo real para ajustar semáforos, desviar flujos y reducir la congestión \texttt{\cite{Aparicio Araque, 2025}}. ¡Menos tiempo en el coche para todos!
  \item \textbf{Servicios Públicos Más Eficientes:} Ayuda a distribuir recursos (ambulancias, bomberos) o a gestionar solicitudes de permisos, ¡haciendo que los servicios sean más rápidos y justos para los ciudadanos!
  \item \textbf{Recomendaciones Personalizadas:} ¿Tu serie favorita en Netflix? ¿Ese producto que `\texttt{justo'' necesitabas en Amazon? Es IA analizando tus gustos en }\texttt{entornos virtuales'' }\cite{Méndez Saldaña, N/A}`.
  \item \textbf{Salud y Diagnóstico:} Asiste a médicos en diagnósticos o en la optimización de tratamientos, siempre buscando el mejor resultado para el paciente, aunque con la sensibilidad de los datos sanitarios y la protección de los mismos \texttt{\cite{Zouak Lara, 2025}}.
\end{itemize}
\textbf{Pregunta Interactiva:} ¿Qué otro ejemplo de IA en su vida diaria se les ocurre?

\textbf{IMG\_PH\_0}

\end{frame}



\begin{frame}[fragile]{¡Sabías que...! Hechos Fascinantes de la IA}


Prepárense para estos datos que les volarán la cabeza:

1.  \textbf{La Mano Invisible de la IA:} ¿Sabías que los algoritmos de IA están tomando cada vez más decisiones que impactan tu vida diaria, desde optimizar el flujo de tráfico y asignar recursos públicos hasta personalizar tus experiencias en línea? A menudo, ni siquiera nos damos cuenta de la influencia de la IA \texttt{\cite{Aparicio Araque, 2025}}. ¡Es la ``mano invisible'' de la IA en acción!
\textbf{IMG\_PH\_0}

2.  \textbf{El Desafío de la `\texttt{Caja Negra'':} Uno de los mayores retos en la toma de decisiones públicas impulsadas por la IA es la \textbf{transparencia}. Imaginen que un sistema de IA deniega una solicitud para un servicio público; explicar \textit{por qué} tomó esa decisión puede ser increíblemente complejo debido a la naturaleza intrincada de los modelos de aprendizaje automático. Este problema de la }\texttt{caja negra'' plantea preguntas significativas sobre la responsabilidad y la equidad cuando la IA se utiliza para el beneficio del público en general }\cite{Zouak Lara, 2025}`.

\textbf{IMG\_PH\_0}

\end{frame}



\begin{frame}[fragile]{Desafíos, Ética y el Futuro: ¿Controlamos la IA?}


Con tanto poder, vienen grandes responsabilidades. La IA y los datos presentan desafíos importantes:

\begin{itemize}
  \item \textbf{Transparencia (La `\texttt{Caja Negra''):} ¿Cómo aseguramos que entendemos \textit{por qué} la IA toma ciertas decisiones, especialmente si afectan nuestra vida? }\cite{Zouak Lara, 2025}`
  \item \textbf{Sesgos y Equidad:} Si la IA aprende de datos históricos, y estos datos tienen sesgos humanos, la IA puede replicar o incluso amplificar esos sesgos. ¿Cómo garantizamos decisiones justas para \textit{todos}?
  \item \textbf{Privacidad y Protección de Datos:} Nuestros datos son valiosos. ¿Cómo se usan, quién tiene acceso y cómo nos aseguramos de que estén protegidos? ¡El \textbf{RGPD} (Reglamento General de Protección de Datos) es un buen ejemplo de cómo se busca protegerlos, incluso frente a la IA! \texttt{\cite{Zouak Lara, 2025}}
  \item \textbf{Rendición de Cuentas:} Si una IA comete un error, ¿quién es responsable? ¿El programador? ¿La empresa? ¿El usuario?
\end{itemize}
\textbf{Pregunta Interactiva:} Si tuvieran que elegir el mayor desafío de la IA en la toma de decisiones, ¿cuál sería y por qué?

\textbf{IMG\_PH\_0}

\end{frame}



\begin{frame}[fragile]{Tu Rol en la Era de la IA: ¡No Seas un Observador!}


La IA no es el futuro, ¡es el presente! Y tú no eres un mero espectador.

\begin{itemize}
  \item \textbf{Sé Informado:} Entiende cómo la IA te afecta. ¡Estás haciendo eso ahora mismo!
  \item \textbf{Haz Preguntas:} Si una decisión automatizada te parece injusta o incomprensible, ¡pregunta por qué! Exige transparencia.
  \item \textbf{Participa:} La regulación y la ética de la IA son temas abiertos. Tu voz como ciudadano es crucial para moldear un futuro donde la IA sea una fuerza para el bien.
  \item \textbf{Empodérate con el Conocimiento:} ¡Aprende! Comprender los conceptos básicos te dará una ventaja enorme.
\end{itemize}
La IA está aquí para quedarse, y su impacto en la toma de decisiones es inmenso. ¡Pero somos nosotros, los humanos, quienes debemos guiarla!

\textbf{IMG\_PH\_0}

\end{frame}



\begin{frame}[fragile]{Conclusión: IA y Datos, Aliados de la Humanidad}


Hemos viajado por el fascinante mundo de la IA y el análisis de datos, descubriendo cómo influyen en nuestras decisiones diarias y en las de la sociedad.

\begin{itemize}
  \item \textbf{La IA es un motor poderoso}, los datos su combustible.
  \item \textbf{Está en todas partes}, desde tu móvil hasta la gestión de ciudades.
  \item \textbf{Tiene un potencial increíble}, pero también grandes desafíos éticos y de transparencia.
  \item \textbf{Tu papel es crucial} para asegurar que la IA se desarrolle de manera responsable y beneficie a todos.
\end{itemize}
¡Gracias por su atención! Espero que se lleven no solo conocimiento, sino también la curiosidad de seguir explorando este emocionante campo.

\textbf{¡El futuro de las decisiones está en nuestras manos... y en las de la IA bien diseñada!}

\textbf{IMG\_PH\_0}

\end{frame}



\begin{frame}[fragile]{Referencias}


\begin{itemize}
  \item Aparicio Araque, Blanca. ``EL IMPACTO DE LA INTELIGENCIA ARTIFICIAL Y LA AUTOMATIZACIÓN DE DATOS EN LA TOMA DE DECISIONES EN EL SECTOR PÚBLICO.'' \textit{Inteligencia y toma de decisiones: perspectivas actuales}, 2025. DOI: 10.69592/979-13-7011-369-8-cap-17.
  \item Méndez Saldaña, M. ``Análisis de datos para la toma de decisiones en entornos virtuales.'' \textit{Tesis Universidad de Concepción}, N/A. DOI: 10.7764/tesisuc/edu/74749.
  \item Zouak Lara, P. ``Sinergias entre el RGPD y la IA en el ámbito sanitario para la toma de decisiones automatizadas.'' \textit{Inteligencia artificial y protección de datos: desafíos en la era digital}, 2025. DOI: 10.69592/979-13-7011-111-3-cap-16.
\end{itemize}
\end{frame}




\begin{frame}[allowframebreaks]{References}
  \bibliographystyle{apalike}
  \bibliography{references}
\end{frame}


\end{document}